\chapter{Discussion} 
%Here you should discuss all aspect of your thesis and project. How did the process work? Which choices did you make, and what did you learn from it? What were the pros and cons? What would you have done differently if you were to undertake the same project over again, both in terms of process and product? What are the societal consequences of your work?

This chapter discusses the results presented in the previous chapters and their implications for the overall system design. The discussion covers the architectural decomposition and timing determinism of the distributed node, the audio-acquisition chain, the analog front end and its \SI{4}{\milli\ampere} IEPE current source, the power-supply noise that proved to be the dominant limitation of the assembled system, the battery autonomy, and the room-acoustic reverberation and sound-insulation measurements, with the objective of evaluating the effectiveness of the proposed distributed measurement system and identifying its limitations.

\section{Architectural Design and System Decomposition}

A key architectural decision in this project was to separate wireless communication and synchronization into a dedicated \acrshort{rf} subsystem. This approach was motivated by the need to achieve predictable, bounded timing behavior in a distributed embedded system, where radio communication introduces asynchronous events and protocol overhead.

By delegating \acrshort{rf} event handling to a dedicated microcontroller and propagating synchronization events through hardware-level signals, the system avoids shared software execution paths that could introduce unbounded latency or jitter. This decomposition enables clear temporal isolation between communication-related processing and time-critical synchronization actions.

Although this architecture increases system complexity compared to monolithic designs, the experimental results demonstrate that the resulting improvement in timing determinism justifies the additional hardware and integration effort.

In addition to its direct effect on timing determinism, the chosen architectural partitioning offers a clear path for system evolution. Since synchronization reaches the audio subsystem as a hardware trigger, rather than via shared software layers or clock domains, the audio processing chain stays independent of any particular wireless protocol or network configuration. This separation allows the RF subsystem to be updated, expanded, or replaced without disturbing the time-critical audio execution path, ensuring the architecture is robust to future changes in communication methods or deployment scale.

A further design decision at the firmware level concerned memory coherency between the CPU cache and the DMA engine. Without dedicated configuration, the cache and the DMA controller would operate on divergent views of the same physical memory, producing intermittent audio corruption that is difficult to reproduce and practically impossible to diagnose at the application level. Marking the shared audio buffers as non-cacheable through the memory protection unit eliminates the hazard entirely, at the cost of slightly increased CPU memory latency, an acceptable trade-off, given that the throughput bottleneck in this system is the SD card I/O rather than cache performance.

\section{4 mA Current-Source Stability and Microphone Headroom}
The 4~mA current source was tested to see how stable the bias current for the IEPE microphone remains as the temperature changes. To do this, the voltage across a 50~$\Omega$ resistor, which acts like a real microphone load, was measured, and Ohm’s law was then used to calculate the current. Ideally, a current of 4~mA should produce 200~mV across the resistor.

At room temperature (about 21$^\circ$C), the current measured 4.032~mA. As the system was cooled down to about $-20^\circ$C, the current increased to 4.383~mA. When it was heated up to around $60^\circ$C, the current dropped to 3.698~mA. This means the current varied by up to +9.6\% or -7.6\% from the ideal value. Over the full $80.6^\circ$C temperature range, the current changed by about $-0.685$~mA, which averages to roughly $-8.5~\mu$A per degree Celsius (about $-0.21\%$ per $^\circ$C). This change was almost perfectly linear; both cooling and heating produced about the same rate of change, indicating that the current’s temperature dependence is predictable rather than abrupt or erratic.

The environmental chamber test used here is much more extreme than what would occur during a typical measurement. In practice, the room temperature changes gradually, and a single measurement takes only a short time. For example, if the temperature changed by 10$^\circ$C, which is much more typical, the current would only change by about 85~$\mu$A, or about 2.1\%. This means that for real-world acoustic measurements, the current source can be considered essentially constant.

The changes in current mostly affect the bias point (operating voltage) for the microphone and the available headroom, the maximum voltage the microphone preamplifier can output without distortion. They do not affect the sensitivity for small signals, which is determined by the microphone capsule and its internal electronics. The important effect is that, if the current gets too low, the preamplifier could run out of headroom and start to distort at high output levels. However, the preamplifier used here (Type~2671) is designed to have a 12~V DC output and handle signals up to 7~V peak (below 20~kHz). With a 24~V supply, there is plenty of headroom (about 5–19~V swing), and the signal remains within specification even if the output voltage varies within the allowed range.

To put these numbers into context, the microphone’s maximum output can be related to sound pressure level (SPL) using the 50~mV/Pa sensitivity of the Type~4189 capsule. A 7~V peak output (4.95~Vrms) translates to about 99~Pa of sound pressure, or roughly 134~dB SPL. This is higher than the designed maximum for the analog front end (AFE), which can accept up to about 3.2~Vrms at the microphone input (about 130~dB SPL). This means that, under normal conditions, the rest of the measurement system will reach its maximum input level before the microphone preamplifier does. Even if the bias current drops to its minimum measured value (a 7.5\% reduction), the loss in headroom is only about $-0.68$~dB, and in a realistic 10$^\circ$C case, only about $-0.18$~dB. These estimates are conservative, since the AFE input is high-impedance and the overall dynamic range is mostly limited by the analog-to-digital converter (codec) and gain settings. Therefore, the effect of current drift on the usable measurement range is small.

From a power perspective, each IEPE channel powered at 24~V uses about 96~mW at the nominal current. Across the full temperature range, the power varies from about 88.8~mW to 105.2~mW. This change is relevant for overall power budgeting and heat generation, but it happens slowly and does not affect the audio signal itself.

In summary, the current source provides nearly 4~mA across a wide temperature range, and its small, predictable drift has little impact on microphone performance. Its influence on headroom is minor compared to the main sources of signal limitation in the system, such as the gain setting and maximum allowable input. This makes it a reliable choice for the IEPE bias in this application.

\section{Power-Supply Regulation, Noise, and System Limitation}
\label{sec:power_supply_noise}
The power-supply results must be analyzed by interpreting three aspects: DC regulation, load-side supply noise, and integration with the audio subsystem. The first aspect verifies whether the designed rails were generated correctly; the second evaluates the residual noise at the filtered outputs; the third explains why a rail that is correct in DC can still degrade the noise floor once connected to a sensitive audio chain, which turns out to be the dominant limitation of the assembled system.

The three output rails reached their intended levels of 3.30003~V, 7.4660~V, and 24.0323~V, with deviations of only about 0.001\%, 0.35\%, and 0.13\% from nominal, respectively. These are small and do not indicate a regulation problem; for the buck converters they fall well within the feedback-network tolerance, whose worst-case ranges (with 1\% resistors) are about 3.25--3.35~V for the 3.3~V rail and 7.31--7.57~V for the 7.4~V rail (implemented with 390~k$\Omega$ and 47~k$\Omega$, the 47~k$\Omega$ substituted for a 50~k$\Omega$ value for availability). The output-terminal ripple was likewise small, 5.33~mV, 0.34~mV, and 0.09~mV, or about 0.16\%, 0.005\%, and 0.0004\% of nominal, confirming stability at the DC level.

The filtering strategy must be evaluated based on the load-side measurements, since the experiment did not isolate each filter stage; thus, the measured spectra include the converter, output filter, PCB routing, load, and probing setup. The selected 4.7~$\mu$H inductors and 22~$\mu$F output capacitors were chosen to stay within the AP63200 operating range and provide current margin (first-order inductor-ripple estimates of about 1.04~A and 1.30~A for the 3.3~V and 7.4~V rails at maximum battery voltage), but these values verify saturation margin rather than predict the operating ripple, particularly under pulse-frequency modulation at light load. Under load, the 24~V rail showed about 27.7~mV peak-to-peak (broad maxima near 9.2~kHz and 17.2~kHz) and the 3.3~V rail about 36.9~mV (maxima near 11.4~kHz and 15.2~kHz), that is, roughly 0.12\% and 1.12\% of the respective rail voltages. The 24~V microphone-bias rail is therefore relatively clean, while the 3.3~V rail is more sensitive in relative terms, as expected for a lower-voltage mixed-domain supply that powers both analog and digital loads.

The implemented filters reduce conducted noise at the supply outputs, but they cannot address all possible noise pathways. Even with output filtering, noise can still enter the system through several mechanisms: shared return impedance (where multiple currents share the same ground path and interfere), inadequate grounding (leading to voltage differences between points that should be at the same potential), magnetic or capacitive coupling from high-frequency switching nodes, suboptimal cable routing, and the supply input of the audio controller board. Therefore, while output filters are essential for noise control, they alone cannot guarantee a clean supply; overall effectiveness depends on careful PCB layout, robust grounding, thoughtful load placement, and well-designed interconnections between circuit boards.

The 7.4~V rail requires a different interpretation, as it powers the STM32H735G-DK via its input path. Integrated measurements revealed that this configuration introduced greater noise than isolated rail tests indicated. This is not due to a DC regulation failure but rather to interactions among the custom power board, the evaluation board input, and the shared ground reference, all of which critically affect the audio noise floor. The prototype’s two-layer in-house PCB, with limited vias and no continuous ground plane, constrained low-impedance returns, tight switching-current loops, and domain separation. 

Additionally, the PCB layout deviated from the AP63200 datasheet recommendations, including using 2~oz copper for lower resistance, adding thermal vias to dissipate heat, implementing a large, continuous ground pour for low-impedance returns, and placing components close together to minimize loop area. These considerations were not taken into account primarily due to fabrication constraints and limited experience with advanced power-supply layout. At a switching frequency of \SI{500}{\kilo\hertz}, high-frequency effects become significant: long or loosely coupled traces, insufficient ground paths, and poor component placement can all lead to increased electromagnetic interference and noise coupling. The lack of these best practices likely resulted in the noise observed when the supply was integrated into the audio chain.

A further, equally important limitation was methodological: the \SI{7.4}{\volt} rail was never exercised under the microcontroller's dynamic load prior to integration. Its bench validation confirmed the correct DC output, but the effectiveness of its \emph{current delivery} under fast load steps was not treated as a design constraint. The implicit assumption was that the evaluation board's on-board regulators would buffer the incoming supply, so the transient current drawn by the custom buck was not characterized. In practice, this proved to be the core problem: the STM32H735G-DK draws short current spikes, peaks on the order of \SI{600}{\milli\ampere} against a typical draw near \SI{390}{\milli\ampere}, and the buck didn’t manage to serve these transients cleanly. As a result, the rail voltage dropped and oscillated whenever the microcontroller suddenly demanded more current, injecting unwanted electrical noise into both the supply input and the ground path. This is a load-transient limitation, distinct from DC accuracy or steady ripple, and would not appear in the unloaded or lightly loaded rail measurements reported above.

The disturbance was attributed to the supply and not to the microphones or the analog conditioning by a control test: physically muting the microphone inputs at the AFE did not clean the recordings while the node was powered from the custom board, whereas powering the same node from a USB power bank or the computer's USB port produced clean channels under otherwise identical conditions. The residual noise thus tracked the power-delivery path, the custom supply together with its unshielded wiring and shared grounding, rather than the acoustic input, isolating that path, and not the microphones or the front end, as the source. This control does not by itself separate the two mechanisms at work in that path: alongside the transient-load coupling described above, the dominant idle component proved to be radiated \SI{50}{\hertz} interference, as the controlled comparison below shows.

These electrical measurements confirm that the rails are correct in DC and carry only modest switching noise at the rail itself, so the limitation is not voltage accuracy but physical integration. The integrated audio measurements make this the most significant finding of the work: the acquisition chain's limiting factor is the power delivery path, not the analog front end or the codec. A controlled comparison conclusively isolated the switching supply and its physical integration. During this test, the STM32 was configured for external power and completely disconnected from the host computer, ensuring the system was fully floating and isolated from the mains grid. When powered by the custom \SI{7}{\volt} supply PCB under these conditions, the idle noise floor rose to \SIrange{-40}{-45}{\decibel FS} and was dominated by a \SI{50}{\hertz} mains-hum comb.

Because the system was running on batteries, the \SI{50}{\hertz} interference was not conducted via a ground loop but was instead picked up as radiated electromagnetic interference (EMI) from the ambient environment. The physical integration relied on unshielded jumper wires to connect the custom power board to the STM32, which created a large loop area. Combined with the custom two-layer PCB's lack of a continuous shielding ground plane, the power delivery path acted as an antenna, coupling ambient \SI{50}{\hertz} room noise directly into the sensitive acquisition chain.

Switching the node's power source to a commercial battery power bank, which eliminated the unshielded jumper wires and utilized a standard, shielded connection, lowered the floor to \SI{-66}{\decibel FS} and completely eradicated the hum. Physically connecting the supply and processor grounds added an additional \SI{5.9}{\decibel} of high-frequency switching noise, confirming that improper grounding is as detrimental to the acquisition chain as inadequate filtering.

Because this disturbance is input-referred, it scales with front-end gain and degrades the high-gain paths: the noise floor rose from \SI{-68.9}{\decibel FS} at unity gain to \SI{-28.3}{\decibel FS} at \SI{+20}{\decibel} and \SI{-8.4}{\decibel FS} at \SI{+40}{\decibel}, each \SI{20}{\decibel} step matching the gain increase. The front end itself is clean, so the high-gain limitation is not due to the conditioning electronics but to the amplification of supply noise. Consequently, the usable dynamic range of the assembled node ($\approx\SI{66}{\decibel}$) was set by the supply, falling short of what the \mbox{16-bit} format and the codec could otherwise provide.

The power-supply board met the voltage-generation requirement. The DC rail deviations were small, the feedback-network behavior was consistent with component tolerances, and the filtered 24~V and 3.3~V rails remained electrically usable under load. The remaining issue is not the nominal voltage accuracy, but the integration of the switching power supply with a sensitive audio measurement chain. The detailed effect of this coupling on the recorded signals is therefore discussed in the acoustic-chain analysis.

\section{Battery Consumption and Autonomy}

The battery subsystem uses a 3S1P lithium-ion pack, which suits the power architecture: it lets the 7.4~V rail be generated with a buck converter and reduces the step-up ratio for the 24~V branch, whereas a 2S pack would fall below the 7.4~V rail during discharge and require a buck--boost stage. It also lowers the battery-side current for a given output power---the 24~V conversion ratio is about 2.16 at the nominal pack voltage versus about 3.24 for a 2S source---which reduces conduction losses in the cells, protection, connectors, and traces.

The analytical autonomy estimate assumed a battery-side base consumption of about 4.81~W (regulated-rail load power with assumed 90\% buck and 85\% isolated-24~V efficiencies). With the 3.35~Ah pack this corresponds to about 37.2~Wh nominal, 29.8~Wh usable at an 80\% factor, and an estimated base autonomy of about 6.2~h---enough for roughly six hours under the assumed base load.

The measured consumption was lower: at about 12.2~V and a mean 0.29933~A, the average battery-side power was about 3.65~W (4.07~W at peak current), roughly 24\% below the analytical estimate. This indicates the design calculation was conservative for the measured operating point, as expected given its assumed efficiencies, nominal currents, and average duty cycle. Using the measured average power as an ideal extrapolation gives about 8.2~h; this is not a validated autonomy result, only a check that the measured current is consistent with the energy budget and that the six-hour target is plausible.

The autonomy validation was limited by the measurement method. The battery current was read with a multimeter rather than a power analyzer or high-resolution logger, so short transients from RF activity, SD-card access, converter mode changes, and playback-dependent loads may not have been captured; the measurement validates the order of magnitude of the consumption but not the full energy profile. As mention before, the base-consumption estimate also excludes the optional active loudspeaker: a passive speaker was connected during the test, which does not represent the worst-case acoustic output, and a high-power playback sequence would raise the average consumption and reduce autonomy.

The system ran from the 3S1P pack for six hours, but the voltage rails were not stable in the long term, and the available data cannot isolate the cause among cell state of charge, internal resistance, protection-circuit behaviour, converter load transients, power-board integration, and measurement limitations. Because of the limited validation time, this subsystem was not characterized as deeply as the synchronization or analog-front-end subsystems.

A simple way to increase autonomy, keeping the same voltage architecture, is a 3S2P pack with the same cells (6.70~Ah, about 74.4~Wh nominal and 59.5~Wh usable), which would give an estimated 12.4~h at the design consumption or 16.3~h at the measured average. This estimation shows that adding capacity in parallel is compatible with the 3S architecture and raises autonomy without changing the regulated rails.

The battery sizing was reasonable for the intended node: the analytical estimate was conservative relative to the measured current, and the 3S configuration suits the required 7.4~V and 24~V rails. However, the validation method was not sufficient to claim a robust operating time; a complete validation would require controlled cell charging, continuous logging of voltage and current, rail monitoring over the full discharge, and representative playback and storage profiles.

\section{Audio Acquisition: Architectural Validation and Determinism}

The audio results provide clear evidence that the main methodology is effective: they show that, when properly configured, distributed embedded hardware can achieve the high timing accuracy required for accurate acoustic metrology measurements. This precision is essential because it makes certain that acoustic data collected across different nodes can be reliably synchronized and compared.

The round-trip latency measurements demonstrated sample-accurate and deterministic behavior; in the electrical domain, the delay was approximately $\SI{39.7}{\ms}$, while acoustically it was $\SI{42.9}{\ms}$ across all repeated tests. Also, the standard deviation was less than one sample, indicating minimal variability among measurements. This consistency confirms that the DMA-driven full-duplex pipeline, SAI clock configuration, and FreeRTOS operating environment introduce a fixed, predictable delay rather than random or variable jitter. Such determinism is essential because it ensures that differences in arrival times among distributed nodes are due to actual acoustic propagation rather than unpredictable timing artifacts in hardware or software.

An important decision regarding development tools was required to achieve this deterministic timing. Initially, the X-CUBE-AUDIO-KIT LiveTuner environment was used to verify hardware functionality and confirm that the signal path was working correctly. However, its always-on streaming approach conflicted with the application’s need for scheduled, SD card-triggered audio acquisition. Recognizing this incompatibility early led to do the proper transition to the BSP audio driver in FreeRTOS for all subsequent development. By making this switch early in the process, the acquisition system was built on a foundation that matched the node’s required deterministic and event-driven behavior, preventing later issues with unpredictable timing.

Although the WM8994 ADC cannot operate at \SI{96}{\kilo\hertz}, even though the DAC is capable of doing so, this limitation does not prevent the system from fulfilling its intended purpose. The audio frequency band captured by the system remains flat within $\pm\SI{0.1}{\decibel}$ from \SI{50}{\hertz} to \SI{19}{\kilo\hertz}, which fully covers the frequency range required for accurate room-acoustic impulse-response and reverberation measurements. Therefore, the ADC’s lower maximum sample rate does not limit the bandwidth needed for the required acoustic analyses.

To evaluate measurement quality, a clear distinction must be made between raw hardware performance and the processing gain of the final acoustic measurements. Raw signal-to-noise ratio (SNR) evaluates the hardware independent of any processing. For this system, the loopback path SNR was measured at \SI{81}{\decibel}, which still represents good performance for practical applications, even though it is lower than the theoretical \SI{98}{\decibel} maximum for \mbox{16-bit} storage.

When the exponential sine-sweep deconvolution method was applied, the recovered impulse achieved an impulse-to-noise ratio (INR) of \SI{125.9}{\decibel} above the pre-arrival noise, far surpassing the physical limits of the \mbox{16-bit} converter. During deconvolution, the energy of the multi-second sweep is coherently integrated into a single, sharp impulse, dramatically boosting the effective signal. In contrast, thermal and quantization noise remains incoherently distributed throughout the recording and is not amplified. 

The decision to use the \mbox{16-bit} sample format was driven by the vendor \acrshort{bsp} library’s default SAI slot width, prioritizing DMA transfer efficiency (Section~\ref{sec:codec_selection}). While a \mbox{16-bit} format theoretically caps dynamic range, the processing gain of the deconvolution pipeline overcomes this limit, making the bit-depth a non-issue for the recovered impulse response. Although upgrading to true \mbox{24-bit} acquisition is a recommended future improvement, it would require significant low-level modifications to the microcontroller's SAI registers and codec configuration. More importantly, until the custom power-supply noise is resolved, the extended dynamic range of a \mbox{24-bit} converter would only digitize the supply-conducted interference with higher fidelity. Once power delivery is improved, \mbox{24-bit} acquisition will allow the system's usable dynamic range to approach the true analog noise floor.

\section{Practical Implications for Acoustic Metrology}

The hardware limitations revealed during testing help explain the differences observed between the two types of room-acoustic measurements performed: reverberation time and sound insulation.

The reverberation-time measurement succeeded because it analyzes the relative decay rate of acoustic energy over time rather than the system's absolute noise floor. In these tests, the energy decay curves displayed a single-slope behavior (with $T_{20} \approx T_{30}$ across all frequency bands of interest). The signal-to-noise ratio was adequate in the mid-frequency range, which allowed the distributed nodes to accurately measure the reverberation characteristics. As a result, the nodes produced consistent and physically meaningful $T_{30}$ values that decreased from approximately $\SI{2.1}{\second}$ to $\SI{0.72}{\second}$, as expected.

A bias of +0.33~\seconds was detected when using the custom estimation routine, as compared to the standard reference routine (\texttt{calcT60\_bands.m}). When the analysis switched to the reference method, this bias disappeared, indicating that the discrepancy originated from methodological choices in the analysis process, not from any hardware limitations. In other words, the difference was due to how the data was processed, rather than any issue with the measurement system itself.

In contrast, the sound-insulation measurement is fundamentally different: it is an absolute dynamic-range assessment that requires the transmitted acoustic signal to exceed the receiving room’s noise floor across all frequency bands. If the signal falls below the noise floor in any band, it becomes impossible to measure the true insulation performance at those frequencies.

The use of a custom power supply led to an increased system noise floor, while the transmitted signal remained modest: the acoustic source was relatively quiet, and the front end operated at unity gain (the \SI{0}{\decibel} follower path) to avoid further amplifying supply noise along with the signal. As a result, at frequencies above approximately $\SI{315}{\hertz}$, the received signal was too weak to rise above the noise floor. Therefore, the measured level difference $D$ and the apparent sound reduction index $R'$ at these higher frequencies could only be reported as lower bounds, since the true signal was masked by noise and could not be accurately determined.

This outcome highlights a limitation in the current experimental setup, rather than a fundamental flaw in the system architecture. The distributed nodes were still able to meet the ISO~16283 measurement margin at low frequencies. With either a louder acoustic source or a lower-noise power supply, the system should be able to achieve full-band sound-insulation measurements.

Beyond the instrumentation, the acoustic measurements introduce their own methodological limitations. While the reverberation and sound-insulation protocols stick to a measurement report prepared by a teaching assistant at NTNU, the absence of on-site guidance from an acoustics specialist meant that certain procedural choices, such as source and microphone placement, spatial averaging, and treatment of low-frequency bands, were subject to operator uncertainty, independent of the hardware used. This uncertainty does not undermine the relative metrics that proved robust, such as the energy-decay slope and low-frequency band-level differences. However, caution is warranted when interpreting the reported absolute insulation levels. Repeating the measurements under the supervision of an experienced acoustician would strengthen confidence in these figures.

Importantly, the distributed wireless nature of the system did not compromise the accuracy of either the reverberation or sound-insulation measurements. This success is due to the inherent nature of these methods: both rely on relative evaluations, such as the energy decay slope or steady-state band-level differences, making them tolerant of network latency. Because these standard acoustic metrology tasks do not rely on exact absolute arrival times across independent nodes, the current wireless synchronization approach proved adequate.

Contrarily, performing distributed impulse-response or time-of-flight measurements, where inter-node delay is used directly to calculate physical distances, will demand much stricter timing. Achieving this requires sample-accurate absolute synchronization across the wireless network, as any timing offset would directly corrupt the spatial measurement. Validating this distributed time base for strict time-of-flight applications is the next critical step in extending the system's capabilities.
\section{Analog Front End (AFE) for IEPE Sensors}

The current-source simulation confirmed that the dual-transistor topology successfully provides about \SI{4.058}{\milli\ampere} of excitation current, well within the Type~2671 specification range. The excitation current ripple is limited to approximately \SI{0.1}{\percent}, demonstrating the high output impedance necessary for proper IEPE microphone operation. This means that while the microphone preamplifier modulates the line voltage to generate the signal, the excitation current remains stable, ensuring accurate sensor operation. 

A key insight from the simulation was the detailed noise analysis, which directly informed the resistor selection process. By understanding how different resistor values contribute to the total noise, as described by thermal (Johnson) noise inherent in all resistors, the design could be optimized to minimize unwanted thermal noise, leading to better overall performance in the finished board.

To make this clear, simulations were run comparing the selected small-valued resistors to much larger alternatives (see Figure~\ref{fig:noise_resistor}). The chosen small resistors ($R_\mathrm{in}=\SI{100}{\ohm}$, $R_f=\SI{1}{\kilo\ohm}$) generate only about \SI{1.8}{\micro\volt} and \SI{0.58}{\micro\volt} of noise, respectively. In contrast, if the resistors were ten times larger, they would contribute around \SI{18}{\micro\volt} and \SI{5.8}{\micro\volt}, a tenfold increase in noise, exactly as predicted by the $\sqrt{R}$ relationship. Overall, this resistor choice reduces the output-referred noise of the gain stage from approximately \SI{22}{\micro\volt} (with large resistors) to about \SI{12}{\micro\volt} (with small ones), as seen in Figure~\ref{fig:noise_opamp}. Therefore, small-valued resistors were intentionally selected to keep the analog conditioning stage’s own thermal noise contribution as low as possible.

This deliberate selection of low-noise resistors is crucial for interpreting the measured noise performance of the finished board. By minimizing the front end’s own thermal noise, any increase in the noise floor at higher gain settings can be traced to external sources, rather than the analog conditioning itself. 

Specifically, as the gain increases, any noise picked up at the input, regardless of its source, is amplified along with the desired signal. For example, in the \SI{0}{\decibel} (follower) setting, there is little gain, so both signal and noise remain low. In higher-gain settings (\SI{20}{\decibel} and \SI{40}{\decibel}), the noise floor rises because all input-referred noise is boosted. 

Detailed measurements showed that this amplified noise was not due to resistor thermal noise, but rather to interference picked up from the power supply. The presence of discrete tones in the output spectrum that align with external supply frequencies rather than the conditioning circuitry, and additional tests using different power sources (see Section~\ref{sec:power_supply_noise}), both confirm the power supply as the main noise contributor. In summary, the resistor noise budget was allocated effectively, and any remaining noise limitation is imposed by the power supply rather than the analog front-end design.

Testing the assembled two-channel board confirmed that all remaining electrical objectives were met. First, the measured gain for each of the four signal paths matched their expected values, within the tolerances set by the resistors used. This means that the amplification provided along each path was as designed. Each gain path clipped only at the predicted input level. In other words, the maximum input that could be processed without distortion was determined by the intended gain structure, not by the limitations of the op-amp’s voltage rails. This is important because it shows that the design limits are well controlled.

The frequency response (passband) remained flat across the entire audio range. This ensures that the system does not introduce unintended coloration or loss of signal fidelity. Harmonic distortion also stayed low as long as each path was operated within its linear range. At unity gain, distortion was about \SI{0.0069}{\percent}, matching the codec-only baseline. At \SI{+20}{\decibel}, distortion ranged from \SI{0.011}{\percent} to \SI{0.014}{\percent}. Distortion increased sharply only after the input signal exceeded the full-scale range. This confirms that the front end does not introduce significant distortion of its own, and that any distortion is determined by the available gain headroom rather than by the electronics' linearity.

Inter-channel isolation through the two-channel PCB measured about \SI{-95}{\decibel} at \SI{1}{\kilo\hertz}. This is roughly \SI{5}{\decibel} below the codec-only baseline of \SI{-100}{\decibel}, but still well below the measurement requirement. This result indicates that the coupling introduced by the board layout is negligible in practice.

In addition to the electrical tests, a brief functional check was performed using a self-built electret microphone. This microphone was assembled on a perforated board and biased by the same \SI{4}{\milli\ampere} current source used in the main design.

Even though the microphone was unshielded and had low sensitivity, it captured the test sweep well enough for deconvolution to recover a single, well-localized impulse response. This was a true end-to-end test, showing that the entire signal path, from microphone through the front end and codec, can transmit a coherent, usable signal. During this test, the idle channel noise stayed below \SI{-70}{\decibel FS}, confirming low background noise when no signal was present.

It is important to note that because of the capsule's low output level and relatively high noise floor, this recording is only suitable for qualitative checking. It demonstrates basic functionality but cannot be used to extract precise quantitative metrics relative to the actual microphones used to run the tests. 

Taken together with the previously demonstrated low intrinsic noise, these results make it clear that the analog front end achieved its design goals. 

\section{Limitations}

The results validate the main design direction of the distributed measurement node, but some aspects of the system were not characterized with the same level of depth. The synchronization subsystem was tested both in the laboratory and outdoors, while the power, battery, and full acoustic integration tests were more limited by available time, instrumentation, and prototype constraints. The limitations discussed in this section define the conditions under which the reported results should be interpreted.

\subsection{RF Range Validation}

The RF range tests validate the feasibility of the selected synchronization link, but they do not fully characterize all field conditions. The integrated-antenna tests were useful to compare \acrshort{los} and terrain-obstructed \acrshort{nlos} behavior, while the external-antenna tests showed that longer distances are possible. However, the outdoor results depend strongly on antenna height, orientation, terrain, vegetation, weather, and node placement.

For this reason, the reported range should not be interpreted as a fixed operating distance of the system. It represents the performance obtained under the tested routes and antenna conditions. A more complete RF characterization would require more repeated measurements, more distance points near the failure region, and stricter control of the antenna setup.

\subsection{Long-Range Synchronization Timing}

The synchronization timing was characterized mainly in the laboratory, where the master and slave outputs could be measured simultaneously with the oscilloscope. This validated the complete local trigger-generation path and showed microsecond-scale alignment between equivalent slave boards.

The long-range outdoor experiments focused mainly on packet success and communication reliability. They did not measure the final trigger alignment at each distance. The laboratory synchronization result should not be directly extrapolated to all outdoor propagation conditions. The outdoor tests verify that synchronization packets can be exchanged at long range, but they do not fully quantify the trigger timing under those same conditions.

\subsection{Power-Supply Noise Isolation}

The power-supply board was validated in terms of DC regulation and load-side noise on the main rails. However, the measurements did not isolate each noise path independently. In particular, the output filters were measured only as part of the complete rail behavior, including the converter, filter, PCB routing, load, and probing setup. Therefore, the attenuation of each filter stage was not quantified separately.

The strongest integration issue appeared through the 7.4~V path used to feed the STM32H735G-DK board, an effect that emerged after system integration rather than as a simple DC regulation failure. Although the control tests in Section~\ref{sec:power_supply_noise} identified radiated \SI{50}{\hertz} interference and load transients as the dominant contributions, the available measurements do not quantitatively separate the individual paths, such as converter ripple, ground coupling, evaluation-board power routing, and cable interconnection.

\subsection{PCB Layout and Grounding Constraints}

The power and interconnection boards were affected by the available fabrication process. The two-layer implementation, limited via options, and restricted ground-plane structure made it difficult to implement ideal low-impedance return paths. This is relevant for switching converters, where high-current loops and shared return impedance can couple noise into sensitive analog or audio paths.

These constraints limit how far the prototype can be used to judge the final achievable noise performance. A board with better ground continuity and cleaner separation between switching, RF, analog, and audio domains could reduce part of the observed coupling.

\subsection{Battery Autonomy Characterization}

The battery measurements verified the approximate order of magnitude of the system consumption, but they did not provide a complete autonomy characterization. The current was measured with a multimeter, so short transients from RF activity, SD-card access, converter mode changes, and playback-dependent loads may not have been captured.

The six-hour battery test should be interpreted only as a functional operating test. The state of charge of the cells, rail stability over time, internal resistance, protection-circuit behavior, and accumulated energy were not monitored with enough resolution. As a result, the battery sizing is supported by the calculations and measured average current, but the long-term autonomy remains only partially validated.

\subsection{Single-Node Acoustic Deployment}
\label{sec:lim_single_node}

Due to time constraints, only one measurement node was fully deployed with the complete analog front end and custom power supply. The reverberation and sound-insulation measurements therefore used an asymmetric two-node configuration: the fully deployed node performed the acquisition, while a second, battery-powered node, with the STM32H735G-DK and CC1352 evaluation board, was supplied from a power bank and served only as the wireless-synchronized source driving the loudspeaker. 

Because only a single analog-front-end PCB was built, a symmetric deployment in which both nodes acquire through their own front ends, capturing the source and receiving rooms at the same instant, could not be realized; the two rooms were instead measured in turn by moving the single acquisition node between them. This is adequate for the steady-state band-level difference used for the insulation index and for the single-room reverberation decay, but it does not yet exploit the simultaneous, time-aligned dual-room capture that the synchronization link is designed to enable.


\section{Broader Implications and Future Work}

The developed system shows that distributed acoustic measurements can be implemented with low-cost embedded nodes when timing, audio acquisition, wireless communication, and power management are treated as coupled design problems. The results also show that the main challenges are not limited to the individual subsystems. RF synchronization, analog signal quality, supply noise, grounding, antenna placement, and battery operation all affect the practical reliability of the final measurement node. Future work should focus on improving the system as an integrated instrument, rather than optimizing each block in isolation.


\subsection{RF Range and Antenna Validation}

The RF tests should be repeated with more controlled antenna conditions. The external-antenna results showed that longer distances are possible, but the achievable range depended strongly on terrain, vegetation, rain, antenna height, orientation, and node placement.

Future tests should document antenna height, polarization, orientation, cable length, and weather conditions at each measurement point. The same route should be tested with integrated and external antennas when possible, using the same packet interval, number of transactions, and node placement strategy.

\subsection{Long-Range Synchronization Timing}

The long-range RF tests should  include timing validation, not only packet success. In this work, synchronization timing was mainly validated in the laboratory, while the outdoor tests focused mostly on communication reliability.

Future measurements should try to record packet success rate, OWT stability, and final trigger alignment at each distance. This would show whether the microsecond-scale synchronization observed in the laboratory is preserved under outdoor long-range conditions.

\subsection{Power-Supply and Audio-Board Integration}

The future work should focus on the integration between the custom power board and the audio-controller board. The isolated DC measurements showed that the voltage rails were generated correctly, but the integrated tests showed that the 7.4~V path used to feed the STM32H735G-DK introduced a stronger noise problem than expected. This path should therefore be redesigned, filtered more carefully, or replaced by a cleaner intermediate supply.

Future tests should measure the 7.4~V rail at different points: directly at the converter output, after the output filter, at the connector, and at the input of the audio-controller board. This would help separate converter ripple, cable or connector coupling, and noise introduced through the evaluation-board power path.

\subsection{PCB Layout and Grounding}

The PCB should be improved by fixing the layout of the power and return paths. The prototype was limited by the available fabrication process, especially the two-layer structure and restricted ground-plane implementation. These constraints likely increased the sensitivity to common-impedance coupling, switching-current loops, and shared return paths.

A future board should prioritize shorter switching loops, better separation between switching, RF, analog, and audio domains, and a lower-impedance ground reference. A multilayer PCB would make this easier, especially for separating noisy converter currents from the analog front end and audio-controller ground.

\subsection{Professional Fabrication and Assembly}

Both prototype boards were fabricated in-house and hand-assembled, which set a practical ceiling on their quality that a future revision should remove by using a standard PCB manufacturer such as JLCPCB. Beyond enabling the multilayer stack-up and continuous ground planes discussed above, professional fabrication would add a proper solder mask over the copper: the prototypes left bare copper exposed on the outer layers, which is harder to solder reliably and leaves unmasked conductor that can couple more readily to nearby nodes, whereas a masked board confines the exposed metal to the intended pads.

Assembly quality was a further limitation. Both boards were soldered by hand by a non-specialist, so joint consistency was variable; as also noted during project supervision, some joints may not fail immediately but could degrade the system's reliability over the longer term. Heat-sensitive components, such as operational amplifiers and buck converters, required particular care during soldering, as excessive heat can damage them. A future build produced by a professional assembly service, or with stencil and reflow assembly, would provide more repeatable joints and reduce this long-term reliability risk.

Finally, a professionally fabricated board provides a complete silkscreen layer that can carry printed labels and operating indicators for the analog front end, identifying the gain-selection paths, connector functions, and channel assignments, so that the node can be operated correctly without relying on external documentation.

\subsection{Output-Filter Characterization}

The output filters should be characterized directly in a future revision. In this work, the rails were measured at the load side, so the measured spectra include the combined behavior of the converter, filter, PCB routing, load, and probing setup.

Future measurements should compare the rail noise before and after each filter stage under the same load condition. This would allow the attenuation of each filter to be quantified and would make it easier to identify whether the remaining noise comes from converter ripple, layout coupling, or the connected load.

\subsection{Battery Autonomy Validation}

The battery subsystem also needs a more complete autonomy validation. The current measurements verified the order of magnitude of the consumption, but they did not fully characterize the energy profile of the node. The use of a multimeter was not enough to capture short transients, converter mode changes, SD-card write peaks, RF bursts, or playback-dependent current variations.

Future tests should use a power analyzer or current logger to record battery voltage, current, power, and accumulated energy during complete operating cycles. The measurements should include idle operation, RF communication, recording, SD-card writing, playback, and speaker-output conditions. A 3S2P battery pack could also be evaluated, since it keeps the same voltage architecture while increasing the available energy.

\subsection{Low-Power Operation}

The autonomy of the node can be improved by reducing the consumption during inactive periods. In the current prototype, the power measurements mainly represent an active or semi-active operating condition. However, in a field measurement scenario, the node may spend a significant amount of time waiting for commands, between acoustic measurements, or storing data without active playback or recording.

Future firmware and hardware revisions should therefore evaluate low-power operating modes. The audio controller could enter sleep or stop modes when acquisition is not active, while the codec, analog front end, 24~V microphone supply, and other unused loads could be disabled between measurements. This would require a defined wake-up sequence using the RTC, GPIO triggers, or scheduled RF activity.

This optimization must be evaluated together with the synchronization requirements. A node that enters a deep sleep mode may reduce autonomy consumption, but it must still wake up early enough to receive synchronization packets and generate the trigger with bounded latency. Future work should characterize the trade-off between standby current, wake-up time, RF availability, and synchronization accuracy.

\subsection{Acoustic Measurement Node Iteration}

For acoustic applications, the next essential hardware improvement is a power-delivery upgrade for the acquisition node. Experimental results demonstrate that the power supply, rather than the analog front end or codec, determines the system’s dynamic range.

A future development should incorporate a low-noise linear post-regulator (LDO) on the analog rail and establish a star-grounded analog reference isolated from switching and digital ground returns. Where practical, battery operation should be preferred, as measurements show it lowers the noise floor by more than \SI{20}{\decibel}.

Additionally, constructing a second analog front-end board would facilitate a symmetric two-node setup, enabling both nodes to acquire signals through their own front ends. This allows simultaneous measurement in both the source and receiving rooms on a shared time base, maximizing the benefits of the synchronization link (Section~\ref{sec:lim_single_node}).


\subsection{Integration Board Fabrication}

A dedicated interconnection board was designed to integrate the audio-controller, analog front end, power supply, and wireless controller. However, it was not fabricated within the project timeline and was therefore excluded from the reported tests.

During measurements, the system was assembled using jumper cables, which introduced uncontrolled parasitic capacitance and inductance, increasing susceptibility to conducted and radiated noise.

Fabricating the integration board is a logical next step. Such a PCB would provide short, controlled interconnections with a defined ground reference, minimizing parasitic effects and noise coupling from jumper wiring.

The integration board should also include input filtering at the development boards' supply inputs. This would attenuate conducted noise from the power supply, identified as the primary limitation of the acquisition chain (Section~\ref{sec:power_supply_noise}), before it reaches the sensitive audio inputs.

